\hypertarget{introduccion}{%
\chapter{Introduccion}\label{introduccion}}

\hypertarget{contexto}{%
\section{Contexto}\label{contexto}}

En los ultimos años, las organizaciones de tecnologia han enfrentado una
presion creciente por acelerar la entrega de software, mejorar la
confiabilidad de sus sistemas y reducir los costos operativos asociados
a infraestructuras tradicionales (Forsgren et al., 2018). Lo que antes
era aceptable, como servidores configurados manualmente, despliegues
coordinados por correo electronico y procesos que dependian del
conocimiento de unas pocas personas, hoy representa un riesgo operativo
y una barrera para la competitividad.

La empresa objeto de este trabajo no es ajena a esta realidad. Durante
años, su infraestructura de aplicaciones se sustento sobre
\textbf{Internet Information Services (IIS)} corriendo en servidores
\textbf{Windows Server on-premises}, con procesos de despliegue en su
mayoria manuales y herramientas de integracion y entrega continua que
fueron quedando obsoletas. Este modelo, si bien funcional en su momento,
presentaba limitaciones cada vez mas evidentes: tiempos de despliegue
prolongados, falta de estandarizacion entre equipos, dificultad para
escalar, y una fuerte dependencia de conocimiento no documentado.

{[}REQUIERE INPUT DEL AUTOR: Descripcion breve de la empresa, de
seguros. Editado para versión final (rubro, tamaño aproximado, cantidad
de equipos de desarrollo) sin datos que la identifiquen directamente.{]}

Frente a este escenario, la organizacion inicio un proceso de
modernizacion que busca transformar la forma en que se construye,
despliega y opera el software. No se trato de un cambio instantaneo ni
de una migracion a la nube publica, sino de una \textbf{transicion
progresiva} hacia una plataforma basada en \textbf{Red Hat OpenShift
on-premises}, incorporando practicas de \textbf{GitOps},
\textbf{infraestructura como codigo} y herramientas modernas de
automatizacion. Complementariamente, se utiliza \textbf{ARO (Azure Red
Hat OpenShift)} para soluciones empaquetadas y servicios especificos que
se benefician de un entorno gestionado en la nube.

\hypertarget{problema}{%
\section{Problema}\label{problema}}

{[}REQUIERE INPUT DEL AUTOR: Confirmar o ajustar la siguiente lista de
problemas. Son los que se asumen en base a las conversaciones iniciales,
pero solo Nicolas puede validar cuales aplican y si falta alguno.{]}

El modelo de infraestructura tradicional de la empresa presentaba una
serie de problemas que impactaban tanto en la operacion diaria como en
la capacidad de evolucion tecnologica:

\begin{itemize}
\tightlist
\item
  \textbf{Despliegues manuales y propensos a error}: los despliegues en
  IIS requerian intervenciones manuales, lo que incrementaba el riesgo
  de errores humanos y hacia que cada puesta en produccion fuera un
  evento de alto estres.
\item
  \textbf{Falta de estandarizacion}: cada equipo o aplicacion podia
  tener un proceso de despliegue diferente, dificultando la operacion,
  el soporte y la incorporacion de nuevos miembros.
\item
  \textbf{Herramientas de CI/CD obsoletas}: las herramientas existentes
  no se adaptaban a las necesidades actuales ni permitian implementar
  practicas modernas de entrega continua. {[}REQUIERE INPUT DEL AUTOR:
  Indicar que herramientas de CI/CD se usaban antes, si se pueden
  mencionar.{]}
\item
  \textbf{Conocimiento concentrado}: los procedimientos dependian del
  conocimiento de pocas personas, generando cuellos de botella y riesgo
  ante rotacion de personal.
\item
  \textbf{Escalabilidad limitada}: escalar aplicaciones o incorporar
  nuevos equipos a la plataforma era un proceso lento y costoso.
\item
  \textbf{Sin versionado de la infraestructura}: no existia un registro
  claro de que estaba desplegado, en que version, ni como reproducir un
  entorno.
\end{itemize}

Cabe aclarar que el ecosistema tecnologico de la empresa no se limita a
IIS. La organizacion tambien gestiona otras plataformas como Java, ECM
(Enterprise Content Management), SQL Server y SAP, ademas de
herramientas de scheduling como Control-M para procesos batch en el area
de operaciones. Sin embargo, el alcance de esta tesis se centra en el
eje principal de la transformacion: \textbf{la migracion de las
aplicaciones desplegadas en IIS hacia una plataforma basada en
contenedores y un modelo moderno de entrega de software}.

Este proceso de modernizacion no solo implico cambios tecnologicos sino
tambien organizacionales. El equipo de administracion de plataformas,
del cual el autor forma parte, fue el impulsor de la transformacion. A
partir de este proceso surgieron nuevos roles que antes no existian en
la empresa, como \textbf{Ingeniero de Plataformas} y \textbf{DevOps},
reflejando la evolucion en la forma de operar y entregar software.
{[}REQUIERE INPUT DEL AUTOR: Confirmar si la descripcion del equipo y
los roles es correcta. Agregar contexto si corresponde: cuando surgieron
estos roles, si hubo otros cambios organizacionales relevantes.{]}

\hypertarget{objetivos}{%
\section{Objetivos}\label{objetivos}}

\hypertarget{objetivo-general}{%
\subsection{Objetivo general}\label{objetivo-general}}

Documentar y analizar el proceso de modernizacion de la infraestructura
de entrega de software de una empresa, desde un modelo basado en IIS
on-premises con despliegues manuales, hacia una plataforma basada en Red
Hat OpenShift con practicas de GitOps, infraestructura como codigo y
automatizacion.

\hypertarget{objetivos-especificos}{%
\subsection{Objetivos especificos}\label{objetivos-especificos}}

\begin{enumerate}
\def\labelenumi{\arabic{enumi}.}
\tightlist
\item
  Describir la situacion inicial de la infraestructura y los procesos de
  la empresa, identificando sus limitaciones.
\item
  Presentar los conceptos teoricos necesarios para comprender las
  tecnologias y practicas adoptadas.
\item
  Documentar la evolucion del proceso de modernizacion a traves de sus
  distintas etapas, desde los primeros comandos manuales hasta la
  implementacion de GitOps con Argo CD.
\item
  Analizar las decisiones tecnicas tomadas en cada etapa y sus
  justificaciones.
\item
  Evaluar los resultados obtenidos, comparando el estado anterior con el
  actual en terminos de eficiencia, estandarizacion y autonomia de los
  equipos.
\end{enumerate}

\hypertarget{alcance}{%
\section{Alcance}\label{alcance}}

La presente tesis cubre los siguientes aspectos:

\textbf{Dentro del alcance:}

\begin{itemize}
\tightlist
\item
  El proceso de adopcion de Red Hat OpenShift on-premises como
  plataforma de contenedores para reemplazar IIS.
\item
  El uso complementario de ARO (Azure Red Hat OpenShift) para soluciones
  empaquetadas (por ejemplo, Keycloak para gestion de identidad).
\item
  La evolucion de los procesos de despliegue: desde comandos manuales,
  pasando por scripts, pipelines, hasta GitOps con Argo CD.
\item
  La implementacion de infraestructura como codigo con Helm y Kustomize.
  {[}PENDIENTE: Helm no esta completamente implementado al momento de
  escritura; se documenta el objetivo y el avance.{]}
\item
  La adopcion del patron App of Apps para la gestion declarativa de
  aplicaciones.
\item
  La incorporacion planificada de herramientas complementarias: Red Hat
  Developer Hub y Ansible.
\item
  El impacto organizacional del cambio: nuevos roles (Ingeniero de
  Plataformas, DevOps), nuevas dinamicas en los equipos.
\end{itemize}

\textbf{Fuera del alcance:}

\begin{itemize}
\tightlist
\item
  La modernizacion de otras plataformas de la empresa (Java, ECM, SAP,
  SQL Server, Control-M).
\item
  Detalles de la arquitectura interna de las aplicaciones
  (microservicios, codigo fuente, etc.).
\item
  Aspectos comerciales o financieros de la transformacion.
\end{itemize}

\hypertarget{metodologia}{%
\section{Metodologia}\label{metodologia}}

El enfoque de esta tesis es \textbf{descriptivo y basado en la
experiencia directa} del autor como parte del equipo que llevo adelante
la modernizacion. Se documenta el proceso tal como ocurrio, con sus
aciertos, errores y decisiones iterativas.

La metodologia incluye:

\begin{itemize}
\tightlist
\item
  \textbf{Observacion participante}: el autor fue parte activa del
  proceso de transformacion.
\item
  \textbf{Documentacion retrospectiva}: se reconstruyen las etapas del
  proceso a partir de documentacion interna, registros de configuracion
  y memoria del equipo.
\item
  \textbf{Comparacion antes/despues}: cada etapa se contrasta con el
  estado previo para evidenciar las mejoras o los desafios introducidos.
\item
  \textbf{Ejemplos tecnicos reales}: se incluyen configuraciones,
  manifiestos y capturas (sanitizados) para ilustrar los conceptos con
  casos concretos.
\end{itemize}

\hypertarget{estructura-del-documento}{%
\section{Estructura del documento}\label{estructura-del-documento}}

El presente documento se organiza de la siguiente manera:

\begin{itemize}
\tightlist
\item
  \textbf{Capitulo 1 - Introduccion} (este capitulo): presenta el
  contexto, el problema, los objetivos y el alcance.
\item
  \textbf{Capitulo 2 - Marco Teorico}: introduce los conceptos
  fundamentales necesarios para comprender las tecnologias y practicas
  adoptadas.
\item
  \textbf{Capitulo 3 - Situacion Inicial}: describe en detalle la
  arquitectura y los procesos de la empresa antes de la modernizacion.
\item
  \textbf{Capitulo 4 - Proceso de Modernizacion}: documenta cada etapa
  de la transicion con detalle tecnico y decisiones tomadas.
\item
  \textbf{Capitulo 5 - Arquitectura Actual}: presenta la plataforma
  resultante y su funcionamiento.
\item
  \textbf{Capitulo 6 - Resultados y Analisis}: compara el antes y el
  despues, evaluando las mejoras logradas.
\item
  \textbf{Capitulo 7 - Conclusiones y Trabajo Futuro}: reflexiones
  finales y lineas de trabajo pendientes.
\end{itemize}
